\documentclass{stex}

\libusepackage{naproche}
\libinput{preamble}

\begin{document}
\begin{smodule}{omega.ftl}

\importmodule[libraries/set-theory]{zero.ftl}
\importmodule[libraries/set-theory]{successor-ordinals.ftl}

\symdef{Infinity}{\omega}

\symdecl*{natural number}
\symdecl*{transfinite induction III}

\begin{definition}[forthel,for={Infinity,natural number}]
  \[ \emph{\Infinity} \DefEq \SClass{n}{\Ordinals}{\classtext{$n \Elem X$ for every $X \Subclass \Ordinals$ such that $\Zero \Elem X$ and for all $x \Elem X$ we have $\Succ(x) \Elem X$}}. \]
  Let a \emph{natural number} stand for an element of $\Infinity$.
\end{definition}

\begin{proposition}[forthel]
  $\Zero \Elem \Infinity$.
\end{proposition}

\begin{proposition}[forthel]
  Let $n \Elem \Infinity$.
  Then $\Succ(n) \Elem \Infinity$.
\end{proposition}

\begin{proposition}[forthel,name=transfinite induction III]
  Let $\Phi \Subclass \Infinity$.
  Assume that $\Zero \Elem \Phi$ and for every $x \Elem \Phi$ we have
  $\Succ(x) \Elem \Phi$.
  Then $\Phi \Eq \Infinity$.
\end{proposition}
\begin{proof}[forthel]
  Suppose $\Phi \NotEq \Infinity$.
  Consider an element $n$ of $\Infinity$ that is not contained in $\Phi$.
  Take $\Phi' \Eq \Phi \Diff \Singleton{n}$.

  (1) $\Zero \Elem \Phi'$.
  Indeed $\Zero \Elem \Phi$ and $\Zero \NotEq n$.

  (2) For each $x \Elem \Phi'$ we have $\Succ(x) \Elem \Phi'$.
  \begin{proof}
    Let $x \Elem \Phi'$.
    Then $\Succ(x) \Elem \Phi$.

    Let us show that $\Succ(x) \NotEq n$.
      Assume $\Succ(x) \Eq n$.
      Then $x \NotElem \Phi$.
      Indeed $n \NotElem \Phi$ and if $x \Elem \Phi$ then
      $n \Eq \Succ(x) \Elem \Phi$.
      Contradiction.
    End.

    Thus $\Succ(x) \Elem \Phi'$.
  \end{proof}

  Therefore every element of $\Infinity$ lies in $\Phi'$.
  Indeed $\Phi' \Subclass \Ordinals$.
  Consequently $n \Elem \Phi'$.
  Contradiction.
\end{proof}

\begin{corollary}[forthel]
  $\Infinity$ is a set.
\end{corollary}
\begin{proof}[forthel]
  Define $f(n) \DefEq \Succ(n)$ for $n \Elem \Infinity$.
  Take a subset $X$ of $\Infinity$ that is inductive regarding $\Zero$ and $f$.
  Indeed $f$ is a map from $\Infinity$ to $\Infinity$.
  Then we have $\Zero \Elem X$ and for each $n \Elem X$ we have $\Succ(n) \Elem X$.
  Thus $X \Eq \Infinity$.
  Therefore $\Infinity$ is a set.
\end{proof}

\begin{proposition}[forthel]
  Let $n \Elem \Infinity$.
  Then $n \Eq \Zero$ or $n \Eq \Succ(m)$ for some $m \Elem \Infinity$.
\end{proposition}
\begin{proof}[forthel]
  Assume the contrary.
  Consider a $k \Elem \Infinity$ such that neither $k \Eq \Zero$ nor $k \Eq \Succ(m)$ for
  some $m \Elem \Infinity$.
  Take a class $\omega'$ such that $\omega' \Eq \Infinity \Diff \Singleton{k}$. %!
  Then $\omega'$ is a set.

  (1) $\Zero \Elem \omega'$.
  Indeed $k \NotEq \Zero$.

  (2) For all $m \Elem \omega'$ we have $\Succ(m) \Elem \omega'$.
  \begin{proof}
    Let $m \Elem \omega'$.
    Then $\Succ(m) \NotEq k$.
    Hence $\Succ(m) \Elem \omega'$.
  \end{proof}

  Thus every element of $\Infinity$ is contained in $\omega'$.
  Therefore $k \Elem \omega'$.
  Contradiction.
\end{proof}

\begin{proposition}[forthel]
  Every element of $\Infinity$ is an ordinal.
\end{proposition}

\end{smodule}
\end{document}
